% ==========================================================
\section{GLSP Architecture \& Framework Selection}
\label{sec:concept_glsp}
% ==========================================================

This section explains the architectural choices and framework selections made for implementing the graphical UVL editor.
It connects those choices to the thesis goal of providing a GLSP-based editor extended for contextual behavioral programs.
It outlines the separation of responsibilities between server and client components, the role of the graphical model used for client-server communication, and the rationale behind selecting the server and client technologies described in the following subsections.

% ----------------------------------------------------------
\subsection{Server Framework Selection: Java}
\label{subsec:glsp_selection_server}
% ----------------------------------------------------------

GLSP provides two production-ready server implementations, one in Java and one in TypeScript.
The Java GLSP server supports multiple source-model strategies.
Besides producing and consuming the native graphical model (GModel) used for client rendering, it integrates directly with Eclipse Modeling Framework (EMF)-based semantic models, optionally retaining both the semantic model and layout data so that visual information can be persisted alongside semantic content.
It can also be configured to use an external EMF repository for centralized persistence, letting the server act as the canonical semantic authority.
The Java server further benefits from a mature ecosystem of multiple examples, reducing implementation friction and providing proven patterns for complex model operations.
The TypeScript server framework, while offering a homogeneous full-stack development experience, lacks the same depth of source model support and is a relatively newer technology platform within the GLSP ecosystem.

The Java server framework is chosen for this work based on the following considerations.
The Universal Variability Language does not provide an existing Ecore metamodel for EMF integration.
However, the UVL creators provide the \texttt{java-fm-metamodel} package, a Java library that parses UVL files using the \texttt{uvl-parser} and materializes them into structured Java objects suitable for programmatic manipulation.
This library enables the GLSP server to load UVL models directly, modify them during editing operations, and re-serialize them back to UVL files, providing a stable foundation for model manipulation.
This is an advantage over the TypeScript server, which would require re-implementing UVL parsing and model manipulation logic from scratch.

Rather than integrating directly with EMF, this work adopts a custom dual-model approach in which the server maintains two parallel representations.
The semantic layer consists of \texttt{java-fm-metamodel} objects that faithfully represent the UVL model and serve as the authoritative source for validation, constraint checking, and model modification.
The visual layer consists of a GModel derived from this semantic representation, which is transmitted to the client and encodes not only model structure but also layout information such as node positions, dimensions, and visual groupings.
Since this layout information is essential for rendering but absent from the UVL source format, the server computes it during the transformation from semantic to visual representation.

By using the GModel as the communication medium, the architecture decouples the semantic representation from its visual counterpart.
The server retains authority over semantic concerns, while the client controls visual rendering, including CSS styling, SVG customization, and interaction handlers.
This division simplifies the client-server protocol and reduces synchronization complexity, as both components operate on well-defined data structures without direct coupling to each other's internal representations.

% RQ2: trade-offs noted for implementation evidence

% ----------------------------------------------------------
\subsection{Client Framework Selection: VS Code}
\label{subsec:glsp_selection_client}
% ----------------------------------------------------------

GLSP provides platform integration packages for multiple integrated development environments, namely Eclipse RCP, Eclipse Theia, and Visual Studio Code (VS Code).
VS Code is selected for client deployment for several reasons.
It is the most widely adopted IDE in contemporary software development practice, reducing adoption barriers for end users.
The platform also supports creating and distributing extensions via the \texttt{.vsix} format, enabling a clean separation between the base UVL tool and the UVL-BP extension as distinct, independently distributed packages.
Finally, the GLSP integration for VS Code is stable, well-documented, and actively maintained, providing a reliable foundation for implementation.
As this work focuses on establishing a proof of concept, a single environment is sufficient; broader IDE support may be explored in future work.

% ----------------------------------------------------------
\subsection{Feature Division: UVL Base vs.\ UVL-BP Extension}
\label{subsec:feature_division}
% ----------------------------------------------------------

A central architectural decision is to separate base UVL functionality from UVL-BP-specific features into two distinct, independently distributed packages.
The base UVL tool is delivered as a standalone \texttt{.vsix} extension providing general-purpose feature modeling capabilities, while the UVL-BP extension builds upon it by adding behavioral programming semantics.
This separation ensures that users who work exclusively with UVL have no dependency on BP-specific tooling, and that the UVL-BP extension can evolve independently without affecting the base editor.

\paragraph{Base UVL tool.}
The base tool renders the feature tree hierarchy with feature names, attributes, parent-child relationships, cardinalities, and constraints.
It displays group relationships like \textit{alternative}, \textit{or} and \textit{group cardinalities} that structure feature relationships and support common cross-tree constraints such as requires and excludes.
Features may be annotated with attributes (arbitrary key-value metadata), and feature-level and group-level cardinalities (e.g., \texttt{[1..3]}) are supported.
Core editing operations include creating and deleting features, modifying relations and constraints, as well as editing attributes.
The graphical representation adheres to the design requirements established in Section~\ref{sec:concept_visualization}.

\paragraph{UVL-BP extension.}
The UVL-BP extension depends on the base UVL tool and introduces additional model elements and runtime capabilities specific to behavioral programming.
It adds two special root features, \texttt{Env} and \texttt{Config}, which represent the context of the model.
A new BThread feature type is introduced to represent behavioral threads as first-class model elements.
Different kinds of behavioral events, namely \textit{requested}, \textit{waited-for}, and \textit{blocked} events, may be attached to a BThread as attributes, encoding the event coordination logic of the behavioral program.
BP-specific constraints as defined by Felber and Götz~\cite{felber_behavioral_2025} are supported through the existing constraint representation of the base tool, extending its semantics without requiring additional model structure.
Finally, the extension realizes a \textit{models@run.time} approach by linking the graphical model to a live behavioral program, reflecting its execution state directly within the editor; this integration is discussed in detail in the following section.